\section{มุมมองรายงานตามบทบาท}
เนื้อหาส่วนนี้อธิบายคำถามที่โครงสร้างข้อมูลและหน้ารายงานรองรับจากการตรวจโค้ด มิใช่ผลทดลองกับผู้ใช้จริง ตัวกรองร่วมประกอบด้วยภาคการศึกษาหรือช่วงเทอม ประเภทลา และช่วงจำนวนวันลา โดยแต่ละหน้ามีตัวกรองรายละเอียดเพิ่มเติม\cite{pages}

\subsection{นิสิต: รูปแบบการลาของตนเอง}
มุมมองนิสิตจำกัดข้อมูลของตนเอง ตอบได้ว่าแต่ละเทอมมีใบลาและวันลากี่รายการ ประเภทใดมีสัดส่วนมาก และรายวิชาใดเกี่ยวข้องกับการลามากที่สุด หน้าเว็บมีภาพรวมและตารางรายเทอม รวมทั้งรายละเอียดรายวิชา การเปรียบเทียบใช้มิติเทอม ประเภทลา และ subject/offering ร่วมกับ fact ทั้งสองระดับ

ข้อมูลนี้ช่วยให้นิสิตทบทวนรูปแบบการเข้าเรียนของตน แต่ไม่เพียงพอสำหรับสรุปอัตราขาดเรียนหรือสิทธิ์เข้าสอบ หากยังไม่มีจำนวนครั้งเรียนทั้งหมดและเกณฑ์ของแต่ละวิชา

\subsection{อาจารย์ที่ปรึกษา: กลุ่มนิสิตในความดูแล}
มุมมองที่ปรึกษาตอบว่าในนิสิตที่ดูแล ชั้นปีใดมีการลาสูง ใครบ้างมีจำนวนใบลาหรือวันลาสูง และรูปแบบเปลี่ยนไปอย่างไรระหว่างเทอม หน้าเว็บมีกราฟเปรียบเทียบชั้นปี สัดส่วน ตารางสรุป Top 10 และการค้นหารหัสนิสิต พร้อมเปิดรายละเอียดรายบุคคล

ขอบเขตกลุ่มนิสิตอ้างอิงความสัมพันธ์ปัจจุบันของที่ปรึกษา ขณะที่มิติที่ปรึกษาใน fact ใช้บริบท ณ เวลายื่น การจัดอันดับที่ปรึกษาในภาพรวมจึงต้องระบุนิยามให้ต่างจากการติดตามรายชื่อนิสิตที่ตนดูแลอยู่ในปัจจุบัน

\subsection{อาจารย์ผู้สอน: ผลกระทบต่อรายวิชา}
มุมมองผู้สอนตอบว่ารายวิชาหรือกลุ่มเรียนที่สอนมีการลาเท่าใด เทอมใดสูงกว่าช่วงอื่น และนิสิตคนใดมีการลาในรายวิชานั้น หน้าเว็บมีกราฟเปรียบเทียบ สัดส่วน ตารางรายวิชา แนวโน้มรายเทอม การกระจายของยอดรายวิชา และรายละเอียดรายวิชากับวันเรียน

หน่วยการวิเคราะห์หลักคือ offering และความสัมพันธ์ใบลากับ offering โดยกรองตามสิทธิ์ผู้สอนปัจจุบัน การมีผู้สอนร่วมต้องไม่เพิ่มยอดรวมรายวิชา และยอดจากหลายวิชาไม่ควรรวมแล้วเรียกว่าใบลารวมของนิสิตโดยไม่ตัดคีย์ซ้ำ

\subsection{หัวหน้าภาควิชา: นิสิตและรายวิชาในภาค}
มุมมองนี้ตอบได้สองแกน ได้แก่ภาพรวมการลาของนิสิตในภาค และผลกระทบต่อรายวิชาที่ภาคเป็นเจ้าของ หน้าเว็บรองรับการเปรียบเทียบตามที่ปรึกษา ชั้นปี ผู้สอน และรายวิชา มีตารางนิสิต การค้นหา และรายละเอียดผลกระทบตามวิชา

แกนนิสิตใช้สังกัดนิสิต ส่วนแกนรายวิชาใช้เจ้าของวิชา นิสิตต่างภาคที่เรียนวิชาของภาคจึงอาจอยู่ในรายงานรายวิชาโดยไม่อยู่ในภาพรวมนิสิตของภาค รายการวิชาในภาคควรรวมวิชายอดศูนย์เมื่อมีบัญชีรายวิชาครบ ไม่เริ่มรายงานจากใบลาที่มีอยู่เพียงอย่างเดียว

\subsection{ผู้บริหาร: ภาพรวมระดับคณะ}
มุมมองผู้บริหารตอบว่าการลากระจายระหว่างภาควิชาอย่างไร หลักสูตรใดมีจำนวนสูง และกลุ่มนิสิตตามที่ปรึกษาใดควรติดตาม หน้าเว็บจัดลำดับภาพรวมตามภาควิชา ต่อด้วยหลักสูตร และที่ปรึกษา พร้อม Top 10 และรายละเอียดภาคภายในคณะที่ได้รับสิทธิ์

จำนวนใบลาเป็นตัวชี้วัดขนาดเหตุการณ์ ยังใช้เปรียบเทียบความเสี่ยงระหว่างภาคโดยตรงไม่ได้ เพราะแต่ละภาคมีจำนวนนิสิตต่างกัน หากต้องการอัตราการลา ต้องเพิ่มตัวหารจากทะเบียนนิสิตที่มีขอบเขตและเทอมตรงกัน การจัดอันดับจึงเป็นจุดเริ่มต้นการตรวจข้อมูล ไม่ใช่ข้อสรุปเชิงสาเหตุ

\subsection{เจ้าหน้าที่: ติดตามภาพรวมภายในภาค}
มุมมองเจ้าหน้าที่มีกราฟเปรียบเทียบตามอาจารย์ สัดส่วน ตารางสรุป และจำนวนการลาแยกชั้นปี ใช้ตอบว่ากลุ่มนิสิตภายใต้ที่ปรึกษาคนใดมีการลามากในช่วงที่เลือก และควรตรวจรายละเอียดกลุ่มใดเพิ่มเติม โดยใช้ขอบเขตภาควิชาของบัญชีตามบริการปัจจุบัน

\subsection{สถานะการเชื่อมข้อมูลที่ตรวจพบ}
บริการฝั่ง API มีโค้ดอ่าน \nolinkurl{dw_*} และตรวจบทบาท แต่ \nolinkurl{LiveReportProvider} ในเว็บเรียก \nolinkurl{createMockLiveReportData()} จึงยังไม่ยืนยันว่าหน้ารายงานเหล่านี้แสดงคลังข้อมูลจริง การมี service และตัวแปลงข้อมูลอยู่แล้วไม่เท่ากับการต่อข้อมูลครบเส้นทาง

นอกจากนี้ ตัวแปลงบางส่วนแปลงค่าชั่วโมงที่ไม่มีข้อมูลเป็นศูนย์ จึงต้องแยกสถานะความครบถ้วนก่อนใช้สรุปว่าไม่มีผลกระทบ บริการอ่านหลายตารางแบบขนานและมีแคชชุดข้อมูล ขณะที่เวลารอบสำเร็จเป็นข้อมูลกำกับ ความสอดคล้องของทุก query ใน snapshot เดียวกันยังต้องตรวจเพิ่มก่อนอ้างว่ารายงานมีความคงที่ระดับรอบนำเข้า
